\section{Linear closure of the cubic seed set}
\label{sec:linear_closure}

The final cubic dictionary is
\[
\mathcal M_3^{\mathcalE}
=
\ClL(\seedM_3^{\mathcalE}).
\]
This section shows explicitly what the closure operation adds and why it
terminates.

\subsection{General action of the linear SEIR field}

Let $q$ be a monomial that contains $b_i$ copies of $E_i$ and $c_i$ copies of
$I_i$. Since $S_i$ is constant under $f_{\mathrm L}$, the product rule gives
\[
\LieL q
=
-\left(\sigma\sum_i b_i+\gamma\sum_i c_i\right)q
+
\sigma\sum_{i:c_i>0}
 c_i\,q\frac{E_i}{I_i}.
\]
The notation $qE_i/I_i$ means that one factor $I_i$ is replaced by $E_i$.
The coefficient $c_i$ records the number of identical infectious factors that
can be differentiated. Thus, for example,
\[
\LieL(S_1I_1^2)
=-2\gamma S_1I_1^2+2\sigma S_1E_1I_1.
\]
Applying the linear derivative again generates $S_1E_1^2$. No new node index
and no new polynomial degree is introduced.

\subsection{Closure chains by monomial family}

The cubic seeds belong to three families. Their closure chains are:

\begin{align*}
SII
&\longrightarrow SIE\longrightarrow SEE,\\
SIE
&\longrightarrow SEE,\\
SSI
&\longrightarrow SSE.
\end{align*}

For example,
\[
S_1I_1I_2
\longrightarrow
\{S_1E_1I_2,\ S_1I_1E_2\}
\longrightarrow
S_1E_1E_2.
\]
Similarly,
\[
S_1S_2I_3\longrightarrow S_1S_2E_3.
\]
The intermediate $SIE$ variants generated from the $SII$ seeds are already
present among the 17 cubic seeds generated directly from the quadratic
exposure observables. Consequently, the only new family types added by closure
in this example are $SEE$ and $SSE$.

\subsection{Coordinates added by closure}

The 12 new $SEE$ monomials are
\[
\begin{aligned}
\{&S_1E_1^2,S_1E_1E_2,S_1E_2^2,\\
  &S_2E_1^2,S_2E_1E_2,S_2E_1E_3,
   S_2E_2^2,S_2E_2E_3,S_2E_3^2,\\
  &S_3E_2^2,S_3E_2E_3,S_3E_3^2\}.
\end{aligned}
\]
The 13 new $SSE$ monomials are
\[
\begin{aligned}
\{&S_1^2E_1,S_1^2E_2,\\
  &S_1S_2E_1,S_1S_2E_2,S_1S_2E_3,\\
  &S_2^2E_1,S_2^2E_2,S_2^2E_3,\\
  &S_2S_3E_1,S_2S_3E_2,S_2S_3E_3,\\
  &S_3^2E_2,S_3^2E_3\}.
\end{aligned}
\]
The family counts are therefore:

% Generated by scripts/generate_example_artifacts.py
\begin{center}
\begin{tabular}{@{}lrrr@{}}
\toprule
Family & Quadratic seeds & Added by closure & Final count\\
\midrule
$SII$ & 12 & 0 & 12 \\
$SIE$ & 17 & 0 & 17 \\
$SEE$ & 0 & 12 & 12 \\
$SSI$ & 13 & 0 & 13 \\
$SSE$ & 0 & 13 & 13 \\
\midrule
Total & 42 & 25 & 67 \\
\bottomrule
\end{tabular}
\end{center}


Hence
\[
|\mathcal M_3^{\mathcalE}|
=
\NumMthreeSeeds+\NumMthreeClosure
=
\NumMthree.
\]

\subsection{Direct closure checks}

Two representative checks are
\[
\LieL(S_1I_1E_1)
=
\sigma S_1E_1^2-(\gamma+\sigma)S_1I_1E_1,
\]
and
\[
\LieL(S_1^2I_1)
=
\sigma S_1^2E_1-\gamma S_1^2I_1.
\]
The newly produced monomials $S_1E_1^2$ and $S_1^2E_1$ are both included in
$\mathcal M_3^{\mathcalE}$.

The accompanying verification script checks the stronger statement
\[
\Mon_3(\LieL q)\subseteq\mathcal M_3^{\mathcalE}
\qquad
\text{for every }q\in\mathcal M_3^{\mathcalE}.
\]
Therefore no degree-three term is omitted when the cubic block is
differentiated.
